% ============================================================
% ChatGIT: Intent-Aware Routing for Multi-Turn Conversational Code Retrieval
% Springer LNCS — 15-page compressed version
% ============================================================

\documentclass[runningheads]{llncs}

\usepackage[T1]{fontenc}
\usepackage[utf8]{inputenc}
\usepackage{graphicx}
\usepackage{booktabs}
\usepackage{amsmath}
\usepackage{amssymb}
\usepackage{multirow}
\usepackage{array}
\usepackage{xspace}
\usepackage[inline]{enumitem}
\usepackage{microtype}
\usepackage[table]{xcolor}
\usepackage{tikz}
\usetikzlibrary{arrows.meta, positioning, backgrounds, calc, fit, patterns}
\usepackage{pgfplots}
\pgfplotsset{compat=1.18}
\usepackage{placeins}
\usepackage{float}
\usepackage{listings}
\lstset{
  basicstyle=\ttfamily\scriptsize,
  keywordstyle=\bfseries\color{cBlue!80!black},
  commentstyle=\color{gray},
  language=Python,
  breaklines=true,
  frame=single,
  rulecolor=\color{gray!40},
  xleftmargin=4pt, xrightmargin=4pt,
  aboveskip=4pt, belowskip=2pt
}
% hyperref loaded last
\usepackage[hidelinks]{hyperref}

\definecolor{cBlue}  {HTML}{3B82F6}
\definecolor{cIndigo}{HTML}{6366F1}
\definecolor{cPurple}{HTML}{A855F7}
\definecolor{cGreen} {HTML}{16A34A}
\definecolor{cDark}  {HTML}{1E293B}
\definecolor{cAmber} {HTML}{D97706}

\newcommand{\chatgit}{\textsc{ChatGIT}\xspace}
\newcommand{\conv}{\textsc{ChatGIT-Bench}\xspace}

\begin{document}

\title{ChatGIT: Intent-Aware Routing for Multi-Turn
Conversational Code Retrieval}

\titlerunning{ChatGIT: Conversational Repository Intelligence}

\author{Anish Gupta\inst{1} \and
        Prakhar Sethi\inst{1} \and
        Ritwik Bhattacharya\inst{1} \and
        Sonia Khetarpaul\inst{1}}

\authorrunning{Gupta, Sethi, Bhattacharya, Khetarpaul}

\institute{
  Computer Science and Engineering,
  School of Engineering and Applied Science,
  Shiv Nadar University (SNU), India\\
  \email{\{ag801, ps385, rb213\}@snu.edu.in,\;
         sonia.khetarpaul@snu.edu.in}
}

\maketitle

\begin{abstract}
Every developer has been lost in an unfamiliar codebase: \texttt{grep} finds
the function but not \emph{why} it breaks, which module owns the logic, or
how the pieces fit together architecturally. \chatgit solves this
conversationally -- paste a GitHub URL, ask in plain English, and receive
cited answers drawn directly from the live source code.

We present \chatgit, a RAG system for multi-turn conversational QA over
GitHub repositories, combining five complementary retrieval novelties through
\emph{intent-aware routing}. A companion web interface exposes an interactive
call-graph visualiser, per-file PageRank/HITS dashboards, a repository tree
explorer, and syntax-highlighted responses with clickable \texttt{file:line}
citations -- giving developers both answers and architectural insight.

Evaluated on \conv (221 turns, 96 conversations, 5 Python repositories),
\chatgit achieves \textbf{Recall@5\,=\,0.832} and \textbf{MRR\,=\,0.594},
outperforming VanillaRAG by $+10.7\%$ R@5 and leading Dense+Reranker on all
metrics. The central finding is a routing effect: naively combining all five
novelties degrades MRR by $-1.6\%$; routing each to its effective intent
yields a $+0.035$ MRR routing gain. Results replicate on five unseen repos.

\keywords{Conversational Code QA \and Retrieval-Augmented Generation
\and Session Memory \and Call Graph \and Intent Classification \and PageRank}
\end{abstract}

% ══════════════════════════════════════════════════════════
\section{Introduction}
\label{sec:intro}

A developer joining a production project faces hundreds of interconnected
files, undocumented design decisions, and years of commit history. Standard
tools -- \texttt{grep}, IDE symbol search, GitHub Code Search -- find tokens
but provide no explanation, no architectural context, and no memory of what
was already looked at. Commercial assistants (Copilot Chat, Cody, Cursor)
retrieve context per-turn with no session state: every follow-up re-fetches
the same entry-point functions, flooding context with already-seen content.

Real developer workflows are \emph{multi-turn} by nature. A debugging session
begins with ``Where is \texttt{authenticate()} called?'', continues with
``Why does it fail when the session is expired?'', and ends with ``Summarise
how the auth module fits together.'' Three fundamental problems emerge:
\textbf{Cross-turn redundancy}: without retrieval-level session state,
every follow-up re-retrieves the same entry-point functions.
\textbf{Structurally-blind retrieval}: cosine similarity cannot distinguish
a core architectural class from a utility used nowhere else, nor stable
foundational code from actively-revised files.
\textbf{Fixed retrieval granularity}: ``Where is \texttt{authenticate()}?''
requires a single tight function chunk; ``Explain the auth module'' requires
file-level summaries. A fixed top-$k$ serves neither.

\chatgit addresses all three through a pipeline of complementary components.
Our contributions are:
\begin{enumerate}[label=(\arabic*),leftmargin=*]
  \item A full multi-turn conversational code QA pipeline combining AST
        chunking, graph-structural retrieval signals, session-aware memory
        with coreference resolution, intent-adaptive granularity, and
        call-graph citation injection (\S\ref{sec:architecture}--\ref{sec:components}).
  \item \conv, a 96-conversation (221-turn) multi-turn retrieval benchmark
        across five Python repositories with per-turn chunk-level ground truth,
        intent annotations, and coreference labels (\S\ref{sec:eval}).
  \item An empirical demonstration that \emph{intent-aware routing} is
        essential: naive combination degrades MRR ($-1.6\%$), while routing
        yields a $+0.035$ MRR routing gain (\S\ref{sec:results}).
  \item A complete ablation showing that each novelty applied naively hurts
        performance (N3: $-0.021$; N4: $-0.012$) but routing yields consistent
        gains on all four intents simultaneously.
\end{enumerate}

% ══════════════════════════════════════════════════════════
\section{Background and Related Work}
\label{sec:related}

\subsection{Retrieval-Augmented Generation and Code Retrieval}

RAG~\cite{lewis2021rag} grounds LLM responses via dense retrieval~\cite{guo2024lightrag}.
BM25~\cite{robertson2009bm25} provides a strong lexical baseline for
code search; CodeBERT~\cite{feng2020codebert} and
GraphCodeBERT~\cite{guo2021graphcodebert} extend BERT to code tokens and
data-flow graphs. RepoCoder~\cite{zhang2023repocoder} iteratively retrieves
for code \emph{completion} (infilling a known function body) -- a task
distinct from our open-ended multi-turn QA setting. RepoAgent~\cite{liu2024repoagent}
generates documentation via LLM + call-graph traversal but is not
interactive. SWE-bench~\cite{jimenez2024swebench} targets automated
issue resolution (code editing), not conversational explanation.

\subsection{Multi-Turn QA and Benchmarks}

CoQA~\cite{reddy2019coqa} and QuAC~\cite{choi2018quac} established
multi-turn QA benchmarks over flat text passages; neither addresses code
repositories, call-graph structure, or retrieval granularity.
Benchmarks CoSQA~\cite{huang2021cosqa} and StaQC~\cite{yao2018staqc}
are single-turn only. ConvCodeWorld~\cite{han2025convcodeworld} benchmarks
multi-turn \emph{code generation} with iterative feedback -- our setting
evaluates multi-turn \emph{retrieval} from full repositories with open-ended
NL questions. \conv fills this gap as the first benchmark for multi-turn,
repository-grounded, intent-annotated retrieval QA.

\subsection{Graph-Based Ranking and Commercial Tools}

PageRank~\cite{brin1998pagerank} has been applied to software dependency
graphs for bug prediction. HITS~\cite{kleinberg1999hits} separates hub
(orchestrator) and authority (core utility) scores. We combine both with
query-conditioned attention and \emph{intent-conditional application} --
the key distinction absent from prior graph-ranking work.
GitHub Copilot Chat, Sourcegraph Cody, Cursor, and Aider all offer
conversational interfaces but expose no retrieval-level session state,
redundancy suppression, or intent-adaptive granularity~\cite{githubcopilot,cody,cursor,aider}.

% ══════════════════════════════════════════════════════════
\section{System Architecture}
\label{sec:architecture}

\chatgit is organised into four online pipeline stages plus a companion web
interface (Fig.~\ref{fig:pipeline}).

\textbf{Stage A: Offline Indexing.} Repositories are cloned, parsed with
language-specific AST analysers (Python, JS/TS, Java, Swift), and converted
to token-aware chunks (512 tokens, 64-token overlap). Chunks are embedded
(BGE-small-en-v1.5) into a ChromaDB persistent index. A function call graph
is built with NetworkX; 500-commit git history derives per-file volatility
scores for N1.

\textbf{Stage B: Query Pre-Processing.} Each user turn undergoes
(i)~intent classification (\textbf{N4}: LOCATE, EXPLAIN, SUMMARIZE, DEBUG),
immediately reconfiguring $k$, rerank depth, and context budget;
(ii)~coreference resolution (\textbf{N3a}: temporal back-references,
repeat-query detection, pronoun resolution) using session history;
(iii)~session summary prepended to the LLM prompt.

\textbf{Stage C: Retrieval and Rescoring.} The expanded query is issued
against the vector index. Candidates are rescored by \textbf{N2} hybrid
PageRank + query-conditioned attention (SUMMARIZE only), \textbf{N1}
git-volatility weights, \textbf{N3b} session-memory penalties (LOCATE/DEBUG),
and a cross-encoder reranker (MiniLM-L-6-v2) with per-file diversity capping.

\textbf{Stage D: Generation.} Selected chunks, session summary, and
\textbf{N5} caller/callee context are assembled into an intent-budget prompt
for Groq Llama-3.1-8B. Post-processing injects \texttt{file:line} citations
and updates session state for subsequent turns.

\textbf{Stage E: Web Interface.} A React + Vite SPA exposes an interactive
call-graph visualiser (vis-network), per-file PageRank/HITS dashboards, a
repository tree explorer, and syntax-highlighted chat with clickable
\texttt{file:line} citations -- giving developers both precise answers and
live architectural insight without leaving the browser.

\begin{figure}[H]
\centering
\resizebox{\linewidth}{!}{%
\begin{tikzpicture}[
  stage/.style 2 args={
    draw=#1!60!black, fill=#1!7, rounded corners=6pt,
    line width=1.0pt, text width=#2, align=left, inner sep=7pt,
  },
  arr/.style={-{Stealth[length=6pt,width=5pt]}, line width=1.2pt, color=#1},
  darr/.style={-{Stealth[length=5pt,width=4pt]}, line width=0.85pt,
    dashed, color=#1},
  arrlbl/.style={draw=#1!50!black, fill=#1!10, rounded corners=2pt,
    font=\scriptsize, inner sep=2pt, text=#1!80!black, align=center},
  every node/.style={font=\small},
]

%% Stage A (full width across top)
\node[stage={cBlue}{11.5cm}] (A) at (0,0) {%
  \textbf{\textcolor{cBlue!80!black}{\large A\enspace Offline Indexing}}
  \hfill\textcolor{cBlue!70!black}{\textit{one-time per repository}}\\[4pt]
  \begin{minipage}[t]{2.5cm}
    \textbullet\ Clone repo\\[2pt]
    \textbullet\ AST parse\\
    \quad\textit{\scriptsize Py/JS/Java/Swift}
  \end{minipage}\hfill
  \begin{minipage}[t]{2.5cm}
    \textbullet\ Chunking\\
    \quad\textit{\scriptsize 512 tok, 64 overlap}\\[2pt]
    \textbullet\ Module summaries
  \end{minipage}\hfill
  \begin{minipage}[t]{2.5cm}
    \textbullet\ BGE-small-en-v1.5\\[2pt]
    \textbullet\ ChromaDB index\\
    \quad\textit{\scriptsize keyed by path}
  \end{minipage}\hfill
  \begin{minipage}[t]{2.5cm}
    \textbullet\ Call graph (NX)\\
    \quad\textit{\scriptsize \textbf{[N5]}}\\[2pt]
    \textbullet\ Git volatility\\
    \quad\textit{\scriptsize 500 commits \textbf{[N1]}}
  \end{minipage}%
};

%% Stage B
\node[stage={cIndigo}{3.0cm}, below=2.2cm of A, xshift=-4.0cm] (B) {%
  \textbf{\textcolor{cIndigo!80!black}{B\enspace Query Prep}}\\[4pt]
  \textbullet\ Receive NL query\\[2pt]
  \textbullet\ \textbf{[N4]} Intent:\\
  \quad LOCATE / EXPLAIN\\
  \quad SUMMARIZE / DEBUG\\
  \quad\textit{\scriptsize sets k, budget}\\[2pt]
  \textbullet\ \textbf{[N3a]} Coref:\\
  \quad\textit{\scriptsize temporal, repeat,}\\
  \quad\textit{\scriptsize pronoun}\\[2pt]
  \textbullet\ Session prepend\\[3pt]
  \textcolor{cIndigo!70!black}{\textit{online, per query}}%
};

%% Stage C
\node[stage={cPurple}{3.0cm}, right=0.55cm of B] (C) {%
  \textbf{\textcolor{cPurple!80!black}{C\enspace Retrieval \&}}\\
  \textbf{\textcolor{cPurple!80!black}{\phantom{C\enspace}Rescoring}}\\[4pt]
  \textbullet\ Cosine vector search\\[2pt]
  \textbullet\ \textbf{[N2]} PageRank\\
  \quad + QCA (SUMM.)\\[2pt]
  \textbullet\ \textbf{[N1]} Volatility wt.\\[2pt]
  \textbullet\ \textbf{[N3b]} Session score\\
  \quad\textit{\scriptsize penalise seen}\\[2pt]
  \textbullet\ MiniLM reranker\\[3pt]
  \textcolor{cPurple!70!black}{\textit{online, per query}}%
};

%% Stage D
\node[stage={cGreen}{3.0cm}, right=0.55cm of C] (D) {%
  \textbf{\textcolor{cGreen!80!black}{D\enspace Generation}}\\[4pt]
  \textbullet\ Intent-budget prompt\\
  \quad\textit{\scriptsize 4K--7K by intent}\\[2pt]
  \textbullet\ \textbf{[N5]} Caller/callee\\
  \quad context inject\\[2pt]
  \textbullet\ Groq Llama-3.1-8B\\[2pt]
  \textbullet\ \texttt{file:line} citations\\[2pt]
  \textbullet\ \textbf{[N3]} Session update\\[3pt]
  \textcolor{cGreen!70!black}{\textit{online, per query}}%
};

%% Arrows
\node[above=0.85cm of B, font=\scriptsize\itshape, text=cIndigo!80!black]
  (querylbl) {``\textit{Where is \texttt{auth()}?}'' -- NL query};
\draw[arr=cIndigo] (querylbl.south) -- (B.north);

\draw[darr=cBlue] (A.south -| C.north) -- (C.north)
  node[near end, right=4pt, arrlbl=cBlue]
    {offline index\\[-1pt](ChromaDB + scores)};

\draw[arr=cIndigo] (B.east) -- (C.west);
\draw[arr=cPurple] (C.east) -- (D.west);

\draw[darr=cGreen] (D.north) -- ++(0,1.5)
  -- ++(-7.6,0) node[midway, above=3pt, arrlbl=cGreen]
    {\textbf{[N3]} session update: seen-chunk IDs $\to$ B}
  -- (B.north);

\draw[arr=cGreen] (D.south) -- ++(0,-0.5)
  node[below, font=\scriptsize\bfseries, text=cGreen!75!black]
    {Answer + \texttt{app/auth.py:47} citations};

\end{tikzpicture}%
}
\caption{\chatgit end-to-end pipeline. Stage~A is a one-time offline build;
Stages~B--D execute per query; Stage~E (web interface) wraps the system.
\textbf{Solid arrows}: per-query data flow. \textbf{Dashed arrows}: pre-built
artifacts and session-memory feedback (D$\to$B). \textbf{[N1--N5]} labels
show which novelty activates at each stage and for which intents.}
\label{fig:pipeline}
\end{figure}

% ══════════════════════════════════════════════════════════
\section{Technical Components}
\label{sec:components}

Table~\ref{tab:novelty_phases} maps each novelty to its pipeline phase and
the intents it is activated for. The routing rule is simple: N3-coref and N4
are \emph{always on} (they never degrade any intent); the other three are
gated. N3-session and N5 hurt when applied to the wrong intent -- N3-session
on SUMMARIZE raises redundancy without improving relevance; N5 on LOCATE
injects caller/callee noise around what should be a pinpoint result. EXPLAIN
is routed to pure cosine with \emph{no} structural boosts -- a positive
design decision, not an omission.

\begin{table}[H]
\centering
\caption{Five \chatgit novelties: pipeline phase and query intents routed to.}
\label{tab:novelty_phases}
\footnotesize
\begin{tabular}{llp{1.8cm}p{1.8cm}}
\toprule
\textbf{N} & \textbf{Phase} & \textbf{Component} & \textbf{Routed to} \\
\midrule
\rowcolor{cBlue!7}
N1  & Indexing   & Git-volatility weighting               & SUMMARIZE \\
\rowcolor{cPurple!7}
N2  & Retrieval  & Hybrid PageRank + QC-attention          & SUMMARIZE \\
\rowcolor{cIndigo!7}
N3a & Query prep & Coreference resolution                  & All intents \\
\rowcolor{cPurple!7}
N3b & Retrieval  & Session-memory scoring + zone bonus     & LOCATE, DEBUG \\
\rowcolor{cIndigo!7}
N4  & Query prep & Intent classification                   & All intents \\
\rowcolor{cGreen!7}
N5  & Generation & Call-graph context injection            & EXPLAIN, DEBUG \\
\bottomrule
\end{tabular}
\end{table}

\begin{center}
\begin{minipage}{0.97\linewidth}
\begin{lstlisting}
def route_components(intent: str) -> set:
    # N3_coref and N4 always active: never hurt any intent
    active = {"N3_coref", "N4"}

    if intent in ("LOCATE", "DEBUG"):
        active.add("N3_session")   # suppress seen chunks

    if intent == "SUMMARIZE":
        active |= {"N1", "N2"}     # PageRank + volatility

    if intent in ("EXPLAIN", "DEBUG"):
        active.add("N5")           # inject callers/callees

    # EXPLAIN -> pure cosine, no structural boosts (by design)
    return active
\end{lstlisting}
\vspace{2pt}
\refstepcounter{figure}%
\noindent{\small\textbf{Fig.~\thefigure.} Intent-aware routing policy.
N3-coref and N4 are always active; remaining components are gated by
query intent. EXPLAIN is explicitly routed to pure cosine.\par}
\end{minipage}
\end{center}

\subsection{Token-Aware AST Chunking and Module Summaries}

Language-specific AST parsers identify function and class boundaries as
primary chunk units (MAX\,=\,512 tokens, 64-token overlap, max 6 chunks
per top-level node). A \emph{module-summary chunk} ($\leq$200 tokens) is
synthesised per file from docstrings and exported symbol signatures.
Module summaries receive $1.20\times$ boost for SUMMARIZE queries; class
definitions receive $1.40\times$ (the primary GT type for architecture
overview questions).

\subsection{Git-History Volatility Weighting (N1)}

A three-factor file volatility score from up to 500 commits:
$\text{vol}(f) = 0.5\,\text{freq}(f) + 0.3\,\text{recency}(f) + 0.2\,\text{diversity}(f)$,
where \textit{freq} is normalised commit count, \textit{recency} is
exponentially-decayed days-since-last-change, and \textit{diversity} is
distinct-committer count. Stable files receive up to $1.40\times$ weight.

The intuition is that heavily revised files are more likely to contain
bugs and design tradeoffs worth discussing, while stable files contain
foundational logic worth emphasising for architecture overview queries.
Routing N1 exclusively to SUMMARIZE avoids degrading LOCATE and EXPLAIN,
where commit history is orthogonal to pinpointing a specific function.
\textbf{Boundary condition:} N1 is file-level; on well-maintained repos
with low inter-file volatility variance it contributes $\pm$0.000 MRR.
Per-function volatility via \texttt{git log -L} is future work.

\subsection{Hybrid PageRank and Graph Attention (N2)}

$h(v,q) = \alpha(q)\cdot\text{PR}(v) + (1{-}\alpha(q))\cdot\text{QCA}(v,q)$
where $\text{QCA}(v,q) = 0.6\,\text{sim}(e_v, e_q) + 0.4\,\overline{\text{sim}}_{u\in\mathcal{N}(v)}(e_u, e_q)$
is query-conditioned neighbourhood attention. Mixture weight $\alpha(q)\in\{0.25, 0.50, 0.70\}$
adapts to query specificity; boost $\text{score}\mathrel{\times}=1+0.05\,h(v,q)$
applied when $h(v,q)>0.3$. HITS~\cite{kleinberg1999hits} hub/authority
scores complement PageRank: hub nodes boosted for SUMMARIZE; authority
nodes for LOCATE/EXPLAIN.

The rationale for routing N2 to SUMMARIZE only is that architecture
overview queries specifically target structurally central classes -- the
ones that appear as hubs in the import graph. For LOCATE queries,
promoting high-PageRank nodes hurts precision because the target function
is often a leaf node unreachable via centrality measures.
\textbf{Boundary condition:} N2 contributes
$-0.002$ MRR overall but is net-positive when routed to SUMMARIZE only.

\subsection{Multi-Turn Session-Aware Memory (N3)}

A per-session \texttt{SessionRetrievalMemory} adjusts scores before reranking:
$s'(c) = s_0(c)\cdot r(c)\cdot z(c)$.
$r(c)$ is the redundancy penalty: $0.30$ (same turn), $0.60$ (one turn ago,
LOCATE), $0.75$ (one turn ago, DEBUG -- softer, allowing legitimate revisit
of error paths), $0.85$ (two+ turns ago). $z(c) = 1 + 0.30\,w_f$ is a
zone-coherence bonus from the session's active-file set, decaying at $0.80^t$
per turn. Module-summary chunks are exempt from redundancy penalties;
SUMMARIZE turns are not recorded in session memory.

Three coreference strategies expand ambiguous queries:
(i)~temporal back-references (``the error from before'');
(ii)~repeat-query detection (identical or near-duplicate turns);
(iii)~pronoun resolution (``it''/``this function'').

The key design insight is that the N3 penalty is applied at the
\emph{retrieval score level}, not post-hoc at the result filter level.
This means the reranker can still recover a penalised chunk if the query
evidence is overwhelming -- the penalty shapes the distribution rather than
hard-blocking chunks. This is why N3 achieves $+0.018$ MRR over ConvAwareRAG
despite both approaches reducing redundancy: N3 penalises, ConvAwareRAG
implicitly reorders but cannot suppress.

\textbf{N3 vs.\ ConvAwareRAG.} ConvAwareRAG appends the previous query
before embedding: it reduces exactly to VanillaRAG at turn~0, \emph{increases}
redundancy on follow-ups (35.2\% vs.\ 28.1\%), and cannot make
intent-conditional decisions. N3's chunk-level scoring achieves $+0.018$
MRR over ConvAwareRAG at comparable overall redundancy.

\subsection{Intent-Driven Granularity Adaptation (N4)}

A trained neural intent classifier (TF-IDF unigrams+bigrams + LinearSVC,
97.4\% 5-fold CV, 604 labeled turns) assigns one of four intents and sets
five retrieval parameters simultaneously (Table~\ref{tab:intent_config}).

The rationale for intent classification is that different developer queries
require fundamentally different retrieval strategies. A LOCATE query
(``Where is \texttt{auth()} defined?'') requires a tight top-1 match with
high precision -- retrieving many chunks dilutes relevance. A SUMMARIZE
query (``How does the auth module fit together?'') requires broad coverage
across all files in a subsystem. A single fixed $k$ cannot serve both well;
N4 lets the retrieval parameters adapt instantaneously to the query intent.

\begin{table}[H]
\centering
\caption{Intent-specific retrieval configuration.}
\label{tab:intent_config}
\footnotesize
\begin{tabular}{lrrrrr}
\toprule
Intent & $k$ & Rerank $n$ & Max/file & Ctx budget & Chunk boost \\
\midrule
LOCATE    & 25 & 6  & 2 & 4\,K  & $1.40\times$ statement \\
EXPLAIN   & 20 & 8  & 3 & 6\,K  & $1.20\times$ function  \\
SUMMARIZE & 12 & 5  & 5 & 7\,K  & $1.40\times$ class\,/\,$1.20\times$ module \\
DEBUG     & 30 & 10 & 4 & 6.5\,K& $1.10\times$ mixed     \\
\bottomrule
\end{tabular}
\end{table}

\subsection{Call-Graph Context Augmentation and Reranking (N5)}

For the top-5 seed chunks with cosine similarity $>0.3$, direct callers
and callees augment the candidate pool with a proportional boost:
$\text{boost}(c_{\text{nb}}) \mathrel{+}= 0.06 \times s_{\text{seed}}$.
Caller/callee chunks are tagged \texttt{[CALLER]}/\texttt{[CALLEE]} in the prompt.
Active for EXPLAIN and DEBUG intents only. A fixed $+0.05$ boost degraded
MRR by $9.7\%$; the proportional formulation proved critical.
After graph rescoring, \texttt{ms-marco-MiniLM-L-6-v2}~\cite{reimers2019sentencebert}
reranks top-$k$ candidates; per-file diversity capping enforces
\texttt{max\_per\_file} from Table~\ref{tab:intent_config}.

The rationale for routing N5 to EXPLAIN and DEBUG only is that call-graph
neighbourhood enrichment improves LLM responses precisely when the user
needs to understand relationships -- why a function fails, or how a module
interconnects. For LOCATE queries the call-graph neighbourhood adds noise:
the caller of \texttt{authenticate()} is irrelevant when the user simply
wants to know the file and line number. The proportional boost (tied to the
seed chunk's own similarity score) ensures that low-confidence seeds do not
flood context with weakly-related neighbours.

\FloatBarrier
% ══════════════════════════════════════════════════════════
\section{Experimental Setup}
\label{sec:eval}

\subsection{\conv\ Dataset Construction}

\conv is constructed in two stages. \textbf{Programmatic generation} parses
the AST of five Python repositories (Flask, Requests, FastAPI, Celery, Click)
and generates templated LOCATE$\to$EXPLAIN$\to$DEBUG conversations (2--4 turns).
\textbf{Ground-truth chunk IDs} follow \texttt{file.py::function\_name};
SUMMARIZE turns additionally include \texttt{file.py::\_\_module\_summary\_\_}.

\conv comprises three tiers: \emph{Tier-1} (56 programmatic multi-turn),
\emph{Tier-2} (34 hand-written multi-turn conversations authored by
developers against real source code), and \emph{Tier-3} (25 single-turn
SUMMARIZE). After filtering 19 conversations lacking evaluable GT chunks,
\textbf{96 conversations (221 turns)} form the evaluation corpus.
A circularity note: Tier-1/3 GT is derived from the same AST chunker.
Mitigation: Tier-2 GT is authored independently; all systems share the
same chunking scheme, making circularity a constant across baselines.

\subsection{Baselines and Metrics}

\textbf{Baselines:} BM25, BM25-SlidingWindow, VanillaRAG (BGE-small-en-v1.5
cosine top-$k$), ConvAwareRAG (previous query appended), HybridRAG
(BM25+BGE linear interpolation~\cite{karpukhin2020dense}), BM25+Reranker,
Dense+Reranker (VanillaRAG top-25 re-ranked by MiniLM-L-6-v2~\cite{nogueira2019passage}).

\textbf{Metrics:} \textbf{Recall@5} is the primary operational metric ---
a missed relevant chunk cannot be recovered downstream, so completeness of
the retrieved set directly predicts answer quality. MRR, NDCG@5, P@1, and
cross-turn redundancy are reported alongside. Statistics: Wilcoxon
signed-rank, bootstrap CI ($n{=}3{,}000$), Bonferroni correction.

\FloatBarrier
% ══════════════════════════════════════════════════════════
\section{Results}
\label{sec:results}

\subsection{Multi-Turn Retrieval on \conv}

Table~\ref{tab:multiturn} shows results (96 conversations, 221 turns).
\chatgit (Routed) achieves \textbf{Recall@5\,=\,0.832} vs.\ 0.354 (BM25)
and 0.752 (VanillaRAG, $+10.7\%$). MRR improves from 0.568 to 0.594
($+4.6\%$, $2.9\times$ over BM25). Against Dense+Reranker: $+9.2\%$ R@5
(0.832 vs.\ 0.763, decisive) and $+6.2\%$ NDCG@5 (0.673 vs.\ 0.634,
decisive); the MRR margin ($+1.3\%$, 0.594 vs.\ 0.586) is within
bootstrap CI. The deployed classifier, \chatgit (Routed+Clf),
achieves MRR\,=\,0.599 and R@5\,=\,0.827, leading all baselines
and confirming that the neural classifier fully closes the oracle/deployed gap.

\begin{table}[H]
\centering
\caption{Multi-turn retrieval on \conv (221 turns, 96 conversations).
\textbf{Bold}: best deployed. CI: bootstrap 95\%.
$\dagger$: oracle GT intent. Redund.: cross-turn chunk repetition.}
\label{tab:multiturn}
\footnotesize
\begin{tabular}{lccccc}
\toprule
System & MRR & R@5 & NDCG@5 & P@1 & Redund. \\
\midrule
\multicolumn{6}{l}{\textit{Lexical baselines}} \\
\rowcolor{gray!8}
BM25 & 0.202$_{\pm.040}$ & 0.354 & 0.220 & 0.081 & 11.7\% \\
BM25-SlidingWindow & 0.194$_{\pm.038}$ & 0.364 & 0.219 & 0.081 & 11.1\% \\
\multicolumn{6}{l}{\textit{Hybrid and reranking baselines}} \\
HybridRAG (BM25+BGE) & 0.492$_{\pm.055}$ & 0.701 & 0.551 & 0.376 & 19.8\% \\
BM25+Reranker & 0.502$_{\pm.061}$ & 0.691 & 0.565 & 0.430 & 19.8\% \\
\multicolumn{6}{l}{\textit{Dense baselines}} \\
VanillaRAG (BGE) & 0.568$_{\pm.057}$ & 0.752 & 0.622 & 0.480 & 28.1\% \\
Dense+Reranker & 0.586$_{\pm.056}$ & 0.762 & 0.634 & 0.489 & 27.6\% \\
ConvAwareRAG & 0.576$_{\pm.056}$ & 0.789 & 0.647 & 0.475 & 35.2\% \\
\midrule
\multicolumn{6}{l}{\textit{ChatGIT ablation}} \\
\chatgit (N3 only) & 0.547$_{\pm.058}$ & 0.775 & 0.625 & 0.457 & 40.8\% \\
\chatgit (N4 only) & 0.556$_{\pm.056}$ & 0.742 & 0.610 & 0.462 & 23.6\% \\
\chatgit (N3+N4) & 0.559$_{\pm.054}$ & 0.802 & 0.642 & 0.457 & 37.2\% \\
\chatgit (all N1--N5) & 0.559$_{\pm.055}$ & 0.802 & 0.644 & 0.462 & 36.5\% \\
\rowcolor{cIndigo!8}
\textbf{\chatgit (Routed+Clf)} & \textbf{0.599}$_{\pm.055}$ & \textbf{0.827}
  & \textbf{0.673} & \textbf{0.502} & 36.3\% \\
\rowcolor{cIndigo!8}
\chatgit (Routed)$^{\dagger}$ & 0.594$_{\pm.055}$ & 0.832 & 0.673 & 0.498 & 36.3\% \\
\bottomrule
\end{tabular}
\vspace{2pt}
{\footnotesize
Routed vs.\ VanillaRAG: $\Delta$R@5\,=\,$+$10.7\%, $\Delta$MRR\,=\,$+$4.6\%.
Routing gain (naive$\to$Routed): $\Delta$MRR\,=\,$+$0.035.
}
\end{table}

\textbf{The routing gap.} Applying all five novelties simultaneously
degrades MRR to 0.559 ($-1.6\%$ vs.\ VanillaRAG, within CI). Routing
each to its effective intent yields $+0.035$ MRR (from 0.559 naive to 0.594
routed, a 6.2-point swing). This demonstrates that \emph{how} components
are combined matters more than \emph{whether} they are included.

\textit{Statistical note.} Overall MRR gain ($+0.026$, $p{<}0.05$) is
borderline; gains on R@5 ($+10.7\%$, $p{<}0.01$) and NDCG@5 ($+8.2\%$,
$p{<}0.05$) are robust. SUMMARIZE: $+34\%$ MRR ($p{<}0.05$). DEBUG:
$+23.6\%$ R@5 ($p{<}0.01$), the strongest per-intent result.

\subsection{Per-Intent and Per-Repository Analysis}

Table~\ref{tab:per_intent} shows per-intent MRR and Recall@5.
\chatgit (Routed) matches or beats VanillaRAG on all four intents.

LOCATE and EXPLAIN achieve \emph{parity at reduced complexity}: LOCATE
matches VanillaRAG exactly (MRR\,=\,0.601, R@5\,=\,0.812); EXPLAIN matches
on MRR (0.624) and improves slightly on R@5 ($+0.014$). This is the correct
routing outcome: routing EXPLAIN to pure cosine confirms that knowing
\emph{when not to apply a component} is as important as applying it correctly.

DEBUG improves $+0.013$ MRR and \textbf{$+23.6\%$ R@5} (0.787 vs.\ 0.637)
via N3 session memory and N5 call-graph context -- the largest per-intent
Recall@5 improvement. SUMMARIZE posts $+34\%$ MRR (0.737 vs.\ 0.548) and
perfect R@5 (1.000 vs.\ 0.712), driven by the $1.40\times$ class-boost
correctly targeting class-definition ground truth.

\begin{table}[H]
\centering
\caption{Per-intent MRR and Recall@5 (221 turns). \chatgit (Routed) matches
or beats VanillaRAG on all four intents.}
\label{tab:per_intent}
\resizebox{\linewidth}{!}{%
\footnotesize
\begin{tabular}{lcccccccc}
\toprule
 & \multicolumn{2}{c}{LOCATE} & \multicolumn{2}{c}{EXPLAIN}
 & \multicolumn{2}{c}{DEBUG} & \multicolumn{2}{c}{SUMMARIZE} \\
\cmidrule(lr){2-3}\cmidrule(lr){4-5}\cmidrule(lr){6-7}\cmidrule(lr){8-9}
System & MRR & R@5 & MRR & R@5 & MRR & R@5 & MRR & R@5 \\
\midrule
BM25 & 0.192 & 0.406 & 0.231 & 0.377 & 0.145 & 0.254 & 0.279 & 0.404 \\
VanillaRAG & 0.601 & 0.812 & 0.624 & 0.813 & 0.477 & 0.637 & 0.548 & 0.712 \\
Dense+Reranker & 0.585 & 0.799 & 0.633 & 0.814 & 0.485 & 0.628 & 0.698 & 0.846 \\
\chatgit (N4 only) & 0.601 & 0.812 & 0.618 & 0.799 & 0.360 & 0.494 & 0.737 & 1.000 \\
\rowcolor{cIndigo!8}
\chatgit (Routed) & \textbf{0.601} & \textbf{0.812}
  & \textbf{0.624} & \textbf{0.827}
  & \textbf{0.490} & \textbf{0.787}
  & \textbf{0.737} & \textbf{1.000} \\
$\Delta$ vs Vanilla & $\pm$0.000 & $\pm$0.000 & $\pm$0.000 & $+$0.014
  & $+$0.013 & $+$0.150 & $+$0.189 & $+$0.288 \\
\bottomrule
\end{tabular}%
}
\end{table}

Per-repository: \chatgit (Routed) improves over VanillaRAG on all five repos.
Flask ($+2.7\%$) and FastAPI ($+7.8\%$) show the largest gains via the
SUMMARIZE class-boost and N5 call-graph context. Click is the only
Dense+Reranker stronghold (0.645 vs.\ 0.597): its flat CLI design means
N5 call-graph neighbours are absent and identifier matching dominates.

\subsection{Generalisation to Unseen Repositories}
\label{sec:extended}

\textbf{Two held-out repos (107 turns).} Tornado (1,905 chunks) and Scrapy
(2,660 chunks) are evaluated with 107 single-turn conversations (programmatic
GT). Three findings replicate: (i)~ConvAwareRAG = VanillaRAG exactly (no
prior context at turn~0 -- turn-0 parity confirmed empirically);
(ii)~All5-naive (0.725) falls below VanillaRAG (0.746), routing recovers
to 0.769 ($+0.044$ routing gain, 25\% larger than on \conv);
(iii)~SUMMARIZE routing generalises ($+14.7\%$ MRR, 0.657 vs.\ 0.573).

\textbf{Large-scale benchmark (275 turns, 5 repos).} An additional 275
single-turn conversations across Django (9,661 chunks), SQLAlchemy (16,306),
Pytest (7,718), Scrapy, and Tornado yield 230 matched turns. The routing
gap replicates: All5-naive (0.628) falls $-2.2\%$ below VanillaRAG (0.642);
routing recovers to 0.644. ConvAwareRAG again equals VanillaRAG exactly.
SUMMARIZE MRR gains: Routed $+2.5\%$, Routed+Clf $+7.3\%$.
Django shows $-14\%$ regression: N4 class boosts amplify high-centrality
framework classes (\texttt{Model}, \texttt{QuerySet}) on large corpora
--- a corpus-size threshold effect addressed in limitations.

\subsection{Latency and LLM-as-Judge Evaluation}

\chatgit adds ${\sim}$2--9\,ms per query over VanillaRAG (intent
classification + session-memory scoring): $\leq$3\,ms on small/medium repos,
8.6\,ms on large FastAPI (2,567 chunks). All systems are well under 50\,ms
retrieval latency.

Llama-3.3-70B (Groq) rated 30 randomly sampled turns on four criteria
(1--5): per-chunk \emph{relevance/completeness} and set-level
\emph{coverage/diversity}. Per-chunk: VanillaRAG 4.33/2.57; \chatgit
4.23/2.37 (near-parity). Set-level: \chatgit achieves \textbf{4.53 diversity
vs.\ 4.27} ($+0.27$), widest on SUMMARIZE (4.89 vs.\ 4.11, $+0.78$).
Coverage near-parity (3.43 vs.\ 3.63). This confirms \chatgit's MRR/R@5
gains arise from set-level diversity improvements, not per-chunk quality
degradation.

\paragraph{Availability.}
\chatgit, \conv, all evaluation scripts, and the trained intent classifier
are publicly available at:
\url{https://github.com/anish-gupta/chatgit}.
The repository includes instructions for loading arbitrary GitHub URLs,
running the full benchmark suite, and reproducing all figures in this paper.

% ══════════════════════════════════════════════════════════
\section{Discussion}
\label{sec:disc}

\paragraph{Intent-aware routing is the central finding.}
The $+0.035$ routing gain (from 0.559 naive to 0.594 routed) exceeds the
absolute gain over VanillaRAG ($+0.026$). Key decisions: (i)~N3 only for
LOCATE/DEBUG; (ii)~pure cosine for EXPLAIN; (iii)~N4 class-boost and N2
PageRank only for SUMMARIZE; (iv)~N5 call-graph only for DEBUG. All four
intents simultaneously match or beat VanillaRAG -- no per-intent trade-offs.
EXPLAIN routing to pure cosine is a positive design decision: knowing
\emph{when not to apply a component} is as important as applying it correctly.

\paragraph{N1/N2 boundary conditions.}
N1 contributes $\pm$0.000 MRR (file-level signal, too coarse for
function-level retrieval; low inter-file volatility variance on well-maintained
repos). N2 contributes $-0.002$ overall but is net-positive when routed to
SUMMARIZE (where entry-point class definitions are the target). Both findings
are transferable: compute per-function volatility for N1; apply PageRank only
for architecture-overview queries for N2.

\paragraph{ConvAwareRAG and N3.}
ConvAwareRAG gains only $+1.3\%$ MRR (within CI) while raising redundancy
to 35.2\%. N3's explicit chunk-level scoring achieves $+0.018$ MRR over
ConvAwareRAG at comparable overall redundancy, through intent-conditional
decisions that concatenation cannot replicate.

\paragraph{Retrieval-grounded coverage.}
CodeBLEU over 108 manually-annotated turns: 0.154 (CI: 0.152--0.157),
stable across all intents (range: 0.147--0.159). ROUGE-L (0.057) and
BERTScore-F1 (0.085) are low due to the raw-code-chunk vs.\ prose-explanation
semantic gap. Uniform CodeBLEU across intents confirms retrieval quality
does not collapse on any intent class.

% ══════════════════════════════════════════════════════════
\section{Limitations}
\label{sec:limitations}

\begin{enumerate}[leftmargin=*,label=(\arabic*)]
  \item \textbf{Scale and CI overlap.} Bootstrap CIs of $\pm$0.055 MRR
        on 221 turns; MRR gain over Dense+Reranker ($+1.3\%$) is within CI.
        Recall@5 ($+9.2\%$, $p{<}0.05$) is decisive. The five repos span
        distinct software domains (web framework, HTTP client, CLI toolkit,
        async framework, task queue) with 2,980--13,035 commits and
        384--9,616 chunks; the benchmark is single-language (Python).
        Scaling to 500+ turns would tighten CIs to $\pm$0.025.

  \item \textbf{N3 redundancy penalty at depth.} In three-turn chains,
        the N3 penalty can suppress the DEBUG ground truth by Turn~2.
        Routing EXPLAIN to pure cosine (no N3) mitigates this for Turn~1,
        but intent-conditional penalty decay at Turn~2+ is future work.

  \item \textbf{N4 class-boost amplification on large corpora.} On Django
        ($>$9K chunks), the $1.40\times$ class boost over-promotes high-centrality
        classes above the GT target ($-14\%$ MRR). Corpus-size-adaptive boost
        scaling is needed.

  \item \textbf{N1/N2 granularity mismatch.} Volatility is file-level;
        retrieval is function-level. Per-function git volatility
        (\texttt{git log -L}) and function-level PageRank are planned.

  \item \textbf{Python-only evaluation.} JS/TS/Java retrieval performance
        is unvalidated. No end-to-end generation evaluation with a live LLM
        has been conducted; a 20-question human Likert rating study is planned.
\end{enumerate}

% ══════════════════════════════════════════════════════════
\section{Conclusion}
\label{sec:conc}

We presented \chatgit, a multi-turn conversational code QA system combining
five complementary retrieval novelties: \textbf{N1} git-volatility weighting,
\textbf{N2} hybrid PageRank + query-conditioned attention, \textbf{N3}
session-aware memory with coreference resolution, \textbf{N4} intent-driven
granularity adaptation, and \textbf{N5} bidirectional call-graph context
injection -- all unified by intent-aware routing.

On \conv (96 conversations, 221 turns, 5 repos), \chatgit (Routed) achieves
Recall@5\,=\,0.832 and MRR\,=\,0.594, outperforming VanillaRAG ($+10.7\%$
R@5, $+4.6\%$ MRR) and leading Dense+Reranker ($+9.2\%$ R@5, $+6.2\%$
NDCG@5). The deployed classifier variant (MRR\,=\,0.599) leads all baselines.

The core finding is \emph{intent-aware routing}: naive combination of all
five novelties degrades MRR ($-1.6\%$); routing yields a $+0.035$ MRR
routing gain with no per-intent trade-offs. SUMMARIZE ($+34\%$ MRR) and
DEBUG ($+23.6\%$ R@5) drive the gains; EXPLAIN is correctly routed to pure
cosine. The routing-gap finding replicates on five additional unseen
repositories. \chatgit, \conv, and all evaluation scripts are available at
\url{https://github.com/anish-gupta/chatgit}.

\paragraph{Acknowledgements.}
This project was implemented by Anish Gupta, Prakhar Sethi, and Ritwik
Bhattacharya under the supervision of Dr.~Sonia Khetarpaul.

% ══════════════════════════════════════════════════════════
\begin{thebibliography}{26}

\bibitem{githubcopilot}
GitHub: GitHub Copilot Chat.
\url{https://docs.github.com/en/copilot/github-copilot-chat} (2024)

\bibitem{cody}
Sourcegraph: Cody --- AI Coding Assistant.
\url{https://sourcegraph.com/cody} (2024)

\bibitem{cursor}
Anysphere Inc.: Cursor --- The AI Code Editor.
\url{https://www.cursor.com} (2024)

\bibitem{aider}
Gauthier, P.: Aider: AI Pair Programming in your Terminal.
\url{https://aider.chat} (2024)

\bibitem{lewis2021rag}
Lewis, P., et~al.: Retrieval-Augmented Generation for
Knowledge-Intensive NLP Tasks. arXiv:2005.11401 (2021)

\bibitem{guo2024lightrag}
Guo, Z., et~al.: LightRAG: Simple and Fast
Retrieval-Augmented Generation. arXiv:2410.05779 (2025)

\bibitem{robertson2009bm25}
Robertson, S., Zaragoza, H.: The Probabilistic Relevance
Framework: BM25 and Beyond.
Found.\ Trends IR \textbf{3}(4), 333--389 (2009)

\bibitem{feng2020codebert}
Feng, Z., et~al.: CodeBERT: A Pre-Trained Model for Programming
and Natural Languages. In: Findings of EMNLP (2020)

\bibitem{guo2021graphcodebert}
Guo, D., et~al.: GraphCodeBERT: Pre-Training Code Representations
with Data Flow. In: ICLR (2021)

\bibitem{zhang2023repocoder}
Zhang, F., et~al.: RepoCoder: Repository-Level Code Completion
Through Iterative Retrieval and Generation. In: NeurIPS (2023)

\bibitem{liu2024repoagent}
Liu, B., et~al.: RepoAgent: An LLM-Powered Framework for
Repository-level Code Documentation Generation.
arXiv:2402.16667 (2024)

\bibitem{kleinberg1999hits}
Kleinberg, J.M.: Authoritative Sources in a Hyperlinked
Environment. J.\ ACM \textbf{46}(5), 604--632 (1999)

\bibitem{brin1998pagerank}
Brin, S., Page, L.: The Anatomy of a Large-Scale Hypertextual
Web Search Engine. Comput.\ Netw.\ ISDN Syst. (1998)

\bibitem{reimers2019sentencebert}
Reimers, N., Gurevych, I.: Sentence-BERT: Sentence Embeddings
using Siamese BERT-Networks. In: EMNLP (2019)

\bibitem{nogueira2019passage}
Nogueira, R., Cho, K.: Passage Re-ranking with BERT.
arXiv:1901.04085 (2019)

\bibitem{reddy2019coqa}
Reddy, S., et~al.: CoQA: A Conversational Question Answering
Challenge. TACL \textbf{7}, 249--266 (2019)

\bibitem{choi2018quac}
Choi, E., et~al.: QuAC: Question Answering in Context.
In: EMNLP (2018)

\bibitem{jimenez2024swebench}
Jimenez, C.E., et~al.: SWE-bench: Can Language Models Resolve
Real-World GitHub Issues? In: ICLR (2024)

\bibitem{huang2021cosqa}
Huang, J., et~al.: CoSQA: 20,000+ Web Queries for Code Search
and Question Answering. In: ACL (2021)

\bibitem{yao2018staqc}
Yao, Z., et~al.: StaQC: A Systematically Mined
Question-Code Dataset from Stack Overflow. In: WWW (2018)

\bibitem{han2025convcodeworld}
Han, H., Hwang, S., Samdani, R., He, Y.: ConvCodeWorld:
Benchmarking Conversational Code Generation in Reproducible
Feedback Environments. In: ICLR (2025)

\bibitem{karpukhin2020dense}
Karpukhin, V., et~al.: Dense Passage Retrieval for Open-Domain
Question Answering. In: EMNLP, pp.\ 6769--6781 (2020)

\bibitem{lin2021few}
Lin, J., Ma, X.: A Few Brief Notes on DeepImpact, COIL, and a
Conceptual Framework for Information Retrieval Techniques.
arXiv:2106.14807 (2021)

\bibitem{petroni2024irrag}
Petroni, F., et~al.: IR-RAG @ SIGIR24: Information Retrieval's
Role in RAG Systems. In: Proc.\ 47th ACM SIGIR (2024)

\bibitem{abootorabi2025multimodal}
Abootorabi, M.M., et~al.: Ask in Any Modality: A Comprehensive
Survey on Multimodal RAG. arXiv:2502.08826 (2025)

\end{thebibliography}

\end{document}
